\documentclass[12pt]{article}


\usepackage[utf8x]{inputenc}
\usepackage{graphicx}
\usepackage[english]{babel}

\usepackage{geometry}
\geometry{hmargin=2.5cm,vmargin=2.5cm}

\usepackage{comment}

\usepackage{xcolor}
\definecolor{color_section}{RGB}{230,8,31}
\definecolor{color_subsection}{RGB}{252,170,50}

\usepackage{hyperref}

\begin{document}
    \begin{titlepage}
        \title{Web technology project plan}
        \author{Mac group}

        \begin{figure}
            \begin{minipage}[c]{0.40\linewidth}
                \includegraphics[width=0.7\textwidth]{TUGraz.png}
            \end{minipage}
            \hfill
            \begin{minipage}[c]{0.48\linewidth}
                \hfill
                \includegraphics[width=0.7\textwidth]{Erasmus.jpeg}
            \end{minipage}
        \end{figure}

        \begin{center}
            \includegraphics[width=0.5\textwidth]{mcdonalds-fat-logo.jpg}\\
            \vspace{1cm}
            \Large McFat Web Application Documentation\\
            \rule{5cm}{.5pt}\\
            \vspace{0.6cm}
            {\Large Web technology}\\
            \vspace{0.1cm}
            {\Large Winter semester 2018-2019} \\
            \vspace{5em}

            \textbf{Jens Mellberg - Matteo Milo$\mathbf{\check{s}}$ - Bastien Pery}\\
        \end{center}

    \end{titlepage}

    \renewcommand{\contentsname}{\color{color_section}\\Table of contents}
    \tableofcontents

    \newpage

    {\color{color_section}\section{Application description}}
    Mac group was composed by four students and four different nationalities.
        Croatian, French, Swedish and Finnish.
        Unfortunately, one of us drop the course just in the middle. 
        But, As part of our Erasmus exchange program, 
        we wanted to develop a web application which could be used by a lot 
        of people around the world and especially by Erasmus people.\\

        \noindent Our application reached this goal because, thanks to it,
        every european student is now able to find important 
        information about McDonald in Europe.\\

    \noindent Our McFat Web Application is divided into four different parts, each including different functionalities:
    \begin{itemize}
        \item MacHome
        \item MacMap
        \item MacStats
        \item MacMenu
    \end{itemize} 

    {\color{color_section}\section{Functionalities}}

        {\color{color_subsection}\subsection{MacHome}}
        \noindent The MacHome page is basically a simple home page which includes 3
        different things. A McDonald logo, a brief description of the
        application and a navigation bar which helps you to reah any other
        part of our application.\\

        {\color{color_subsection}\subsection{MacMap}}
        \noindent The MacMap page is a little bit more complex. Indeed, like
        all other pages, it includes the navigation bar but also a map and one
        button.\\

        \noindent The map is rendered by an external map API provided by 
        \url{https://www.esri.com}. Thanks to this API, you can easily
        navigate on the map, zoom-in or zoom-out on it. Also, you can see 
        that almost all european McDonald locations are displayed on it.
        (We linked some data with the API). If you click on one location,
        it will display the latitude, the longitude and the city of this
        location.\\

        \noindent The second thing on this page is the button "Display your
        coordinates and find the nearest McDonald!". If you click on this button
        our application will give you your location (if you accept it) and give
        you how far is the nearest McDonald and in which city!

        \newpage

        {\color{color_subsection}\subsection{MacStats}}
        \noindent On the MacStats page you can find some new different elements.
        Indeed, on this page, thanks to 2 different forms, users can display
        some statistics about European McDonalds.\\

        \noindent For example, you can select as many countries as you wish
        and choose which data you want to know such as number of restaurants
        in these countries, number of restaurants per 100k inhabitants... When
        you click on the button "Draw  Statistic", some charts will appear and
        show you the data you asked for.\\
        {\color{color_subsection}\subsection{MacMenu}}

    {\color{color_section}\section{Web technologies used}}

    \noindent In our application we used a lot of different web technologies.\\

    \noindent First of all AJAX. Our application is an asynchronous web application.
    To perform this task, we used the jQuery API.\\

    \noindent We also used some HTML5 technologies which were required:\\
    \begin{itemize}
        \item Canvas as graphic technology for the statistic charts on MacStats
        \item WebStorage as database to store all the McDonald locations and do some analysis on it
        \item WebWorkers on the MacMenu page which allow you to send your menu by mail to someone
        \item The Geolocation API to find the user's location on MacMap
        \item Some HTML5 forms on MacStats
        \item The drag and drop HTML5 techonlogy on MacMenu.
    \end{itemize}


    {\color{color_section}\section{Separation of concerns}}

    {\color{color_section}\section{Context of implementation}}

\end{document}